\worksheettitle
  {Чтение и запись многозначных чисел}
  {\modulemark{M04} Учебный модуль \textbullet\ базовый маршрут и запас практики}

\studentnameblock

\begin{checkpointbox}[Вход: готовность после M03]
  Разделите число на классы и подпишите их:
  \[
    50804007=\answerblank[55mm].
  \]
  Группа класса тысяч: \answerblank[18mm]. Число единиц этого класса:
  \answerblank[18mm].
\end{checkpointbox}

\section{Как читать число по классам}

Рассмотрите запись:
\[
  \underbrace{214}_{\text{миллионы}}\;|
  \underbrace{036}_{\text{тысячи}}\;|
  \underbrace{508}_{\text{единицы}}.
\]

Двигайтесь слева направо: назовите число каждого ненулевого класса, затем
название класса. У класса единиц название не произносят.

\begin{fillbox}[Читаем по шагам]
  \begin{enumerate}[label=\textbf{\arabic*.}]
    \item (214) миллионов --- «двести четырнадцать миллионов».
    \item (036) тысяч содержит число \answerblank[15mm] ---
      «\answerblank[50mm]».
    \item (508) единиц --- «\answerblank[55mm]».
  \end{enumerate}
  Всё число: \answerblank[125mm].
\end{fillbox}

\begin{notebox}[Нулевой класс]
  В числе (73\,000\,406) класс тысяч записан группой (000). При чтении
  его не называют: «семьдесят три миллиона четыреста шесть».
\end{notebox}

\section{Тысяча, тысячи или тысяч?}

Соедините число с подходящим словом:
\[
  21\;\answerblank[22mm]\qquad
  32\;\answerblank[22mm]\qquad
  15\;\answerblank[22mm]\qquad
  104\;\answerblank[22mm].
\]

Слова для выбора: \emph{тысяча, тысячи, тысяч}.

\begin{practicebox}[\levelbase. Прочитайте числа]
  Запишите каждое число словами.
  \begin{enumerate}[label=\textbf{\arabic*.}]
    \item $6\,018$ --- \answerblank[105mm].
    \item $42\,305$ --- \answerblank[100mm].
    \item $700\,040$ --- \answerblank[95mm].
    \item $9\,006\,012$ --- \answerblank[85mm].
    \item $305\,000\,080$ --- \answerblank[75mm].
  \end{enumerate}
\end{practicebox}

\newpage
\section{Как записывать число цифрами}

Запишем: \emph{сорок два миллиона семь тысяч девять}.

\begin{center}
\renewcommand{\arraystretch}{1.55}
\begin{tabular}{c|c|c}
  \toprule
  миллионы & тысячи & единицы \\
  \midrule
  $42$ & $007$ & $009$ \\
  \bottomrule
\end{tabular}
\end{center}

Получаем $42\,007\,009$. В непервом классе всегда три позиции, поэтому
$7$ записали как $007$, а $9$ --- как $009$.

\begin{reflectionbox}[Алгоритм записи]
  \begin{enumerate}[label=\textbf{\arabic*.}]
    \item Выделите названные классы.
    \item Запишите число каждого класса отдельной группой.
    \item В каждом непервом классе дополните группу \answerblank[20mm]
      нулями до \answerblank[12mm] цифр.
    \item Неназванный промежуточный класс запишите как \answerblank[18mm].
    \item Соедините группы и прочитайте результат обратно.
  \end{enumerate}
\end{reflectionbox}

\begin{practicebox}[\levelbase. Запишите цифрами]
  Сначала при необходимости подпишите группы классов.
  \begin{enumerate}[label=\textbf{\arabic*.}]
    \item двадцать восемь тысяч пять: \answerblank[55mm];
    \item шестьсот три тысячи сорок два: \answerblank[55mm];
    \item семь миллионов девятьсот тысяч восемь: \answerblank[55mm];
    \item пятьдесят миллионов двадцать четыре: \answerblank[55mm];
    \item двести один миллион три тысячи семьдесят: \answerblank[55mm].
  \end{enumerate}
\end{practicebox}

\begin{warningbox}[Исправьте запись]
  Число \emph{восемь миллионов пять тысяч сорок два} записали так:
  $8\,5\,42$.

  Правильные группы: \answerblank[20mm]\(\;|\;\)\answerblank[20mm]
  \(\;|\;\)\answerblank[20mm].
  Число: \answerblank[42mm].

  Объяснение ошибки: \answerblank[90mm].
\end{warningbox}

\section{Закрепление: выбирает учитель}

Эта часть не обязательна для перехода. Выполните столько заданий, сколько
нужно для уверенности.

\begin{practicebox}[\levelpractice. Чтение]
  Запишите словами.
  \begin{enumerate}[label=\textbf{\arabic*.}]
    \item $18\,004$; \item $90\,090$; \item $407\,001$;
    \item $12\,030\,004$; \item $80\,005\,600$;
    \item $906\,000\,015$.
  \end{enumerate}
  \vspace{18mm}
\end{practicebox}

\begin{practicebox}[\levelpractice. Запись]
  Запишите цифрами и проверьте обратным чтением.
  \begin{enumerate}[label=\textbf{\arabic*.}]
    \item сто пять тысяч четырнадцать;
    \item девятьсот тысяч девять;
    \item два миллиона сорок тысяч сто;
    \item семьдесят миллионов шесть тысяч два;
    \item четыреста миллионов четыре;
    \item триста два миллиона пятьсот тысяч пятьдесят.
  \end{enumerate}
  \vspace{10mm}
\end{practicebox}

\begin{fillbox}[Пары: взаимная проверка]
  Один ученик записывает число словами, второй --- цифрами. Затем поменяйтесь
  ролями и проверьте друг друга обратным действием.
  \begin{enumerate}[label=\textbf{\arabic*.}]
    \item $61\,008\,030$;
    \item сто семь миллионов двадцать тысяч шесть.
  \end{enumerate}
\end{fillbox}

\section{Разбираемся в ошибках}

\begin{warningbox}[Найдите и объясните ошибки]
  Для каждого пункта запишите правильный ответ и причину ошибки.
  \begin{enumerate}[label=\textbf{\arabic*.}]
    \item «Пять миллионов двенадцать» записали $5\,012$.
    \item $34\,000\,080$ прочитали: «тридцать четыре миллиона восемьдесят
      тысяч».
    \item «Двести тысяч семь» записали $200\,7$.
    \item $7\,014\,005$ прочитали: «семь миллионов четырнадцать тысяч
      пятьсот».
  \end{enumerate}
  \vspace{30mm}
\end{warningbox}

\newpage
\begin{checkpointbox}[Контрольная точка]
  Выполните без подсказки.
  \begin{enumerate}[label=\textbf{\arabic*.}]
    \item Прочитайте $63\,000\,407$:
      \answerblank[75mm].
    \item Запишите цифрами: \emph{девяносто миллионов восемь тысяч
      пятнадцать}: \answerblank[28mm].
    \item Почему во втором числе нужны группы $008$ и $015$?
      \answerblank[102mm].
  \end{enumerate}
\end{checkpointbox}

\begin{reflectionbox}[Моя готовность]
  \choicebox\ Читаю число по классам и пропускаю только нулевой класс.

  \choicebox\ При записи сохраняю все классы и три позиции внутри группы.

  \choicebox\ Мне нужна таблица классов M04-R.
\end{reflectionbox}

\textbf{Закончили раньше?} Выполните M04\(\star\): найдите две разные ошибки
и составьте собственную ловушку.
